\chapter{Smoothers: Jacobi, Chebyshev, and Hiptmair}
\label{ch:smoothers}

% local math macro: eigenvector (not in simbook.sty)
\providecommand{\vphi}{\boldsymbol{\varphi}}

\chapterorient{The previous chapter argued that a good preconditioner is an
approximate inverse cheap enough to apply at every iteration. This chapter builds
the cheapest useful one. A \emph{smoother} is a relaxation so simple it makes a
laughable standalone solver---left to run on its own it crawls---yet it has a
hidden talent the eigenvalue picture brings out at once: it annihilates the
\emph{oscillatory} part of the error in a sweep or two while barely touching the
\emph{smooth} part. That lopsidedness is not a defect. It is exactly the division
of labor a multigrid method needs, and the smoother is the workhorse bolted into
every level of the V-cycle we build in \cref{ch:multigrid}. We meet three: damped
Jacobi (a diagonal scaling), Chebyshev (a fixed-degree polynomial tuned to a
spectral window), and Hiptmair's distributive relaxation (the one fix that makes
edge-element curl--curl problems tractable at all). Each is, in the end, an
apply-function of the kind \cref{ch:precond} taught us to plug in.}

\section{Error in the eigenvector basis}

Everything in this chapter follows from one change of coordinates. Suppose we are
solving a symmetric positive-definite system
\begin{equation}
  \mat{A}\vx = \vb,
  \label{eq:smoo-system}
\end{equation}
the kind that comes out of every weak form in \cref{part:fieldtheory}. Let
$\hat\vx$ be the exact solution, and let an iterative method produce a running
guess $\vx^{(m)}$ after $m$ steps. The object that actually matters is not the
guess but its \emph{error},
\begin{equation}
  \ve^{(m)} = \vx^{(m)} - \hat\vx.
  \label{eq:smoo-error}
\end{equation}
A solver is good exactly when it drives $\ve^{(m)}$ to zero quickly. To see
\emph{how} an iteration acts on the error, we write the error in the natural
coordinate system of $\mat{A}$: its eigenvectors.

Because $\mat{A}$ is symmetric positive-definite, it has a full set of orthonormal
eigenvectors $\vphi_1,\dots,\vphi_n$ with positive eigenvalues
\begin{equation}
  \mat{A}\vphi_k = \lambda_k\,\vphi_k,
  \qquad 0 < \lambda_1 \le \lambda_2 \le \dots \le \lambda_n,
  \label{eq:smoo-eig}
\end{equation}
and we order them from smallest to largest. Any error vector expands uniquely in
this basis,
\begin{equation}
  \ve^{(m)} = \sum_{k=1}^{n} c_k^{(m)}\,\vphi_k,
  \qquad c_k^{(m)} = \inner{\ve^{(m)}}{\vphi_k}.
  \label{eq:smoo-expand}
\end{equation}
The coefficients $c_k^{(m)}$ are the error's \emph{spectral content}: $c_k^{(m)}$
measures how much of the current error lies along the $k$th eigenvector. Driving
the error to zero \emph{means} driving every $c_k$ to zero.

\begin{notationbox}
We use ``frequency'' as a shorthand for ``eigenvalue index $k$.'' For the model
problems in this chapter the eigenvectors $\vphi_k$ really are sinusoids, and the
index $k$ is a genuine spatial frequency: small $k$ is a slowly-varying, smooth
mode; large $k$ is a rapidly-oscillating one. For a general $\mat{A}$ the
eigenvectors are not sinusoids, but the intuition survives---\emph{low frequency}
means the smooth, near-null-space end of the spectrum (small $\lambda_k$), and
\emph{high frequency} the oscillatory, large-$\lambda_k$ end. The whole subject of
smoothing is the story of what an iteration does, eigenvalue by eigenvalue, to the
coefficients $c_k$.
\end{notationbox}

The reason this coordinate change is so powerful is that the simplest iterations
act on each coefficient \emph{independently}. A relaxation built from $\mat{A}$ and
a multiple of the identity (or of the diagonal) commutes with $\mat{A}$'s
eigenvectors, so in the eigenvector basis it is \emph{diagonal}: it multiplies each
$c_k$ by a number that depends only on $\lambda_k$. Call that number the
iteration's \emph{damping factor} at frequency $k$,
\begin{equation}
  c_k^{(m+1)} = g(\lambda_k)\, c_k^{(m)},
  \qquad\text{so}\qquad
  c_k^{(m)} = \bigl[g(\lambda_k)\bigr]^{m}\, c_k^{(0)}.
  \label{eq:smoo-damping}
\end{equation}
The function $g$ is the \emph{amplification factor} or \emph{smoothing symbol} of
the iteration. A component shrinks fast where $\abs{g(\lambda_k)}$ is small and
slowly where $\abs{g(\lambda_k)}$ is near $1$. The entire behavior of a smoother is
encoded in the single curve $g$ plotted against $\lambda$. \Cref{fig:smoo-errorfreq}
shows what this looks like on an actual error: a sum of a smooth mode and an
oscillatory mode, with the oscillatory part visibly gone after a couple of sweeps
and the smooth part stubbornly intact.

\begin{figure}[t]
\centering
\begin{tikzpicture}
\begin{axis}[
  width=12.6cm, height=5.4cm,
  xlabel={position along the domain},
  ylabel={error},
  xmin=0, xmax=1, ymin=-1.35, ymax=1.6,
  xtick=\empty, ytick={-1,0,1},
  legend style={font=\scriptsize, at={(0.5,1.02)}, anchor=south,
    legend columns=3, draw=none, /tikz/every even column/.append style={column sep=8pt}},
  axis lines=left, tick style={simGray}, every axis label/.style={font=\small},
  samples=200, domain=0:1,
]
  % smooth + oscillatory error (initial)
  \addplot[simRust, very thick] {0.85*sin(180*x) + 0.55*sin(180*9*x)};
  % after a few sweeps: oscillatory damped, smooth survives
  \addplot[simBlue, thick, densely dashed] {0.80*sin(180*x) + 0.06*sin(180*9*x)};
  % the smooth part alone (target the smoother leaves behind)
  \addplot[simGray, thin, dotted] {0.80*sin(180*x)};
  \legend{initial error, after 3 sweeps, smooth part}
\end{axis}
\end{tikzpicture}
\caption[Smoothing removes the oscillatory error]{An error built from a smooth
mode ($k=1$) and an oscillatory mode ($k=9$). After three damped-Jacobi sweeps
(dashed) the oscillatory ripples are essentially gone, but the smooth underlying
shape (dotted) is almost untouched. A smoother is precisely the iteration that
does this: it is a high-pass error filter. The smooth remainder is the part the
coarse grid of \cref{ch:multigrid} will handle.}
\label{fig:smoo-errorfreq}
\end{figure}

\begin{keyidea}
A smoother is an iteration whose damping factor $g(\lambda)$ is small at the
\emph{high-frequency} (large-$\lambda$) end of the spectrum and near $1$ at the
\emph{low-frequency} (small-$\lambda$) end. It therefore kills oscillatory error
in a handful of cheap sweeps but barely moves smooth error---a \emph{terrible}
standalone solver and the \emph{essential} ingredient of multigrid, which delegates
the leftover smooth error to a coarse grid. ``Smoother,'' not ``solver,'' is the
right word: its job is to make the error smooth, not small.
\end{keyidea}

\section{Damped Jacobi: diagonal smoothing}
\label{sec:smoo-jacobi}

The simplest relaxation imaginable takes a step proportional to the residual,
preconditioned by nothing more than the diagonal of $\mat{A}$. Write
$\mat{D} = \operatorname{diag}(\mat{A})$ for the matrix holding $\mat{A}$'s
diagonal entries and zeros elsewhere. \emph{Damped Jacobi} is the update
\begin{equation}
  \vx \;\leftarrow\; \vx + \omega\,\mat{D}^{-1}\bigl(\vb - \mat{A}\vx\bigr),
  \label{eq:smoo-jac-update}
\end{equation}
where $\vr = \vb - \mat{A}\vx$ is the residual and $\omega>0$ is a \emph{damping
factor} we are free to choose. Each step costs exactly one operator apply
$\mat{A}\vx$ (\cref{ch:precond}'s matrix-free matvec) plus an elementwise scaling
by $\mat{D}^{-1}$---no inner products, no triangular solves, nothing global. It is
the cheapest nontrivial iteration in the book.

\subsection{The diagonal inverse as a constructed operator}

In the synthesis surface the operator $\omega\mat{D}^{-1}$ is not a matrix; it is a
vector applied elementwise. Building it is a two-step pipeline: \emph{extract} the
diagonal of $\mat{A}$ into a vector, then take its reciprocal entry by entry.
Palace constructs the smoother once, at setup, and stores only the result.

\begin{lstlisting}[style=l4style]
-- setup: build the damped inverse-diagonal smoother once
jacobi_setup :: (A: LinOp[N,N], omega: Float) -> JacobiSmoother[N]
jacobi_setup A omega =
  let d    = assemble_diagonal A      -- extract diag(A) as a vector
      dinv = reciprocal d             -- elementwise 1 / diag(A)
  in  JacobiSmoother { dinv = scale omega dinv }   -- absorb omega into dinv

-- apply: a single elementwise product. no loop, no reduction.
jacobi_smoother :: (op: JacobiSmoother[N], x: Tensor[N]) -> Tensor[N]
jacobi_smoother op x = op.dinv `hadamard` x
\end{lstlisting}

Two things are worth pausing on. First, the damping factor $\omega$ is
\emph{absorbed} into the stored vector \lstinline[style=l4style]|dinv| at setup
time, so the per-call action is a bare elementwise product---``the thinnest
constructed-operator gate'' in the whole framework. Second, the apply is
\emph{inner-product-free}: it touches each entry once and never forms a global sum.
On a parallel machine it requires no communication at all, a property it shares
with Chebyshev and that will matter enormously when we distribute the work.

\begin{notationbox}
\lstinline[style=l4style]|assemble_diagonal| reads the operator's main diagonal;
for a matrix-free high-order operator it returns an \emph{approximate} diagonal,
and the smoother tolerates that---an approximate $\mat{D}^{-1}$ is still a valid
linear map, only a slightly weaker preconditioner. The diagonal extraction, the
reciprocal, and the elementwise product are three firm primitives; the Jacobi
smoother is their composition and nothing more.
\end{notationbox}

\subsection{The damping factor per frequency}

Now run \cref{eq:smoo-jac-update} on the error. Subtract the exact solution
$\hat\vx$ (which satisfies $\mat{A}\hat\vx=\vb$) from both sides. Using
$\vb = \mat{A}\hat\vx$, the residual becomes $\vr = -\mat{A}\ve$, and the update on
the error is
\begin{equation}
  \ve \;\leftarrow\; \ve - \omega\,\mat{D}^{-1}\mat{A}\,\ve
       \;=\; \bigl(\mat{I} - \omega\,\mat{D}^{-1}\mat{A}\bigr)\ve.
  \label{eq:smoo-jac-iter}
\end{equation}
The error is multiplied by the \emph{iteration matrix}
$\mat{G} = \mat{I} - \omega\mat{D}^{-1}\mat{A}$. To read off the per-frequency
damping we diagonalize. For a model problem the diagonal is a constant multiple of
the identity, $\mat{D} = \delta\mat{I}$, so $\mat{D}^{-1}\mat{A}$ shares $\mat{A}$'s
eigenvectors with eigenvalues $\lambda_k/\delta$. Writing $\mu_k = \lambda_k/\delta$
for the scaled eigenvalues---which run over an interval $[\mu_{\min},\mu_{\max}]$
---the damping factor at frequency $k$ is
\begin{equation}
  \boxed{\;
  g(\mu_k) = 1 - \omega\,\mu_k.
  \;}
  \label{eq:smoo-jac-symbol}
\end{equation}
This single affine function is the entire content of damped Jacobi. It is large
(near $1$) when $\mu_k$ is small and decreases as $\mu_k$ grows. \emph{High
frequencies are damped harder than low frequencies}---which is the smoothing
property we wanted---but only if $\omega$ is chosen well.

\begin{pitfall}
Undamped Jacobi ($\omega = 1$) can \emph{fail to smooth at all}, and can even
diverge. With $\omega=1$ the symbol is $g(\mu)=1-\mu$. If the largest scaled
eigenvalue exceeds $2$, then $g(\mu_{\max}) < -1$ and the highest-frequency error
\emph{grows}. Even when it does not diverge, the symbol at the top of the spectrum
($\mu \approx \mu_{\max}$) can sit near $g=-1$, so the most oscillatory modes are
barely damped---the opposite of what a smoother must do. Damping is not a tuning
nicety; without it Jacobi is not a smoother.
\end{pitfall}

\subsection{Choosing the damping factor}

We want $\omega$ to make $\abs{g}$ as small as possible across the
\emph{high-frequency half} of the spectrum, the band a multigrid coarse grid cannot
resolve. The classical balance equalizes the symbol at the two ends of the target
band $[\mu_{\mathrm{lo}}, \mu_{\max}]$:
\begin{equation}
  g(\mu_{\mathrm{lo}}) = -\,g(\mu_{\max})
  \quad\Longrightarrow\quad
  1 - \omega\mu_{\mathrm{lo}} = -(1-\omega\mu_{\max})
  \quad\Longrightarrow\quad
  \omega^\star = \frac{2}{\mu_{\mathrm{lo}} + \mu_{\max}}.
  \label{eq:smoo-jac-optomega}
\end{equation}
This is exactly the formula Palace uses on its estimated-damping path: it estimates
the top eigenvalue $\mu_{\max}$ of $\mat{D}^{-1}\mat{A}$ by a few power-iteration
steps and sets $\omega^\star = 2/(\mu_{\min} + \mu_{\max})$. The spectral interval
$[\mu_{\min},\mu_{\max}]$ comes from the estimate; no exact eigenvalues are ever
formed. The optimal-$\omega$ formula is the standard Jacobi convergence
result~\citep[\S4.1]{saad2003}.

\begin{workedexample}{Damped Jacobi on a 1-D model problem}{smoo-jacobi}
Take the canonical test: the 1-D Laplacian on $n$ interior points with unit
spacing, the tridiagonal matrix $\mat{A}=\operatorname{tridiag}(-1,2,-1)$. Its
eigenvectors are the discrete sinusoids $\vphi_k$ with
$(\vphi_k)_j = \sin\!\bigl(\tfrac{k j\pi}{n+1}\bigr)$ and eigenvalues
\[
  \lambda_k = 2 - 2\cos\!\Bigl(\frac{k\pi}{n+1}\Bigr)
            = 4\sin^2\!\Bigl(\frac{k\pi}{2(n+1)}\Bigr),
  \qquad k = 1,\dots,n.
\]
The diagonal is $\mat{D}=2\mat{I}$, so $\mu_k = \lambda_k/2$ ranges over roughly
$(0,2)$: the lowest mode $k=1$ has $\mu_1\approx 0$ and the highest mode $k=n$ has
$\mu_n\approx 2$. The half-spectrum we want to damp is $[\,1,\,2\,]$ (the modes the
coarse grid misses). Balancing the ends with \cref{eq:smoo-jac-optomega} gives
\[
  \omega^\star = \frac{2}{1 + 2} = \frac{2}{3}.
\]
Now compare the damping at a representative low frequency and a representative high
frequency. At $\mu \approx 0.05$ (a smooth mode) the symbol is
$g = 1 - \tfrac{2}{3}(0.05) \approx 0.967$: after one sweep that mode retains
$96.7\%$ of its amplitude. At $\mu = 2$ (the most oscillatory mode)
$g = 1 - \tfrac{2}{3}(2) = -\tfrac{1}{3}$: that mode retains only $33\%$, a sign
flip and a two-thirds reduction. After $3$ sweeps the high mode is down to
$\abs{g}^3 = (1/3)^3 \approx 0.037$---a $27\times$ reduction---while the smooth mode
is still at $0.967^3 \approx 0.90$. The smoother has cut the oscillatory error by a
factor of about $25$ relative to the smooth error in three cheap sweeps. That ratio
is the whole point.
\end{workedexample}

\Cref{fig:smoo-jacsymbol} plots the symbol $\abs{g}$ for three choices of $\omega$
and shows graphically what the worked example computed: with $\omega^\star=2/3$ the
curve is pinned small across the entire high-frequency half, while undamped
$\omega=1$ leaves the top of the spectrum essentially undamped.

\begin{figure}[t]
\centering
\begin{tikzpicture}
\begin{axis}[
  width=12.6cm, height=6.0cm,
  xlabel={scaled eigenvalue $\mu = \lambda/\delta$ \ (low freq.\ $\to$ high freq.)},
  ylabel={damping factor $|g(\mu)|$},
  xmin=0, xmax=2, ymin=0, ymax=1.05,
  xtick={0,0.5,1,1.5,2}, ytick={0,0.25,0.5,0.75,1},
  legend style={font=\scriptsize, at={(0.02,0.04)}, anchor=south west, draw=simGray!50},
  axis lines=left, tick style={simGray},
  samples=200, domain=0:2,
]
  % shaded high-frequency target band [1,2]
  \addplot[draw=none, fill=simBlue!10, forget plot] coordinates {(1,0)(2,0)(2,1.05)(1,1.05)} \closedcycle;
  \node[font=\scriptsize, simBlue, anchor=north] at (axis cs:1.5,1.02) {high-freq.\ band};
  % undamped omega = 1: |1 - mu|
  \addplot[simRust, very thick] {abs(1 - x)};
  % omega = 2/3 (optimal for [1,2])
  \addplot[simBlue, very thick] {abs(1 - 0.6667*x)};
  % omega = 1/2 (over-damped)
  \addplot[simGray, thick, densely dashed] {abs(1 - 0.5*x)};
  \legend{$\omega=1$ (undamped), $\omega=2/3$ (optimal), $\omega=1/2$}
\end{axis}
\end{tikzpicture}
\caption[The Jacobi damping factor versus frequency]{The damping factor
$\abs{g(\mu)}=\abs{1-\omega\mu}$ for damped Jacobi. Undamped Jacobi ($\omega=1$,
rust) returns to $1$ at the top of the spectrum---it does not smooth the highest
modes. The optimal $\omega=2/3$ (blue) holds $\abs{g}$ below $1/3$ across the entire
high-frequency band $[1,2]$, exactly where the multigrid coarse grid needs it small.
No choice of $\omega$ makes $\abs{g}$ small at low frequency ($\mu\to0$, where
$g\to1$); that is structural, and it is why a smoother alone cannot solve.}
\label{fig:smoo-jacsymbol}
\end{figure}

The figure makes the structural limitation visible: \emph{every} curve returns to
$g=1$ as $\mu\to0$. No diagonal relaxation can damp the smooth modes, because at
$\mu=0$ the update \cref{eq:smoo-jac-update} does nothing. This is not a failure of
tuning. It is the reason smoothers are paired with coarse grids, and it motivates
the next two sections, which sharpen the high-frequency damping without ever trying
to fix the low.

\section{Chebyshev smoothing: a polynomial tuned to the spectrum}
\label{sec:smoo-cheby}

Damped Jacobi takes one step of the affine symbol $1-\omega\mu$. We can do far
better by taking, in effect, several coordinated steps---applying a
\emph{polynomial} in $\mat{D}^{-1}\mat{A}$ whose graph is engineered to be small
across the whole high-frequency band at once. After $k$ such coordinated steps the
error is
\begin{equation}
  \ve \;\leftarrow\; p_k\!\bigl(\mat{D}^{-1}\mat{A}\bigr)\,\ve,
  \label{eq:smoo-cheby-action}
\end{equation}
where $p_k$ is a polynomial of degree $k$ with $p_k(0)=1$ (a unit constant term, so
that the iteration does nothing to the zero-residual fixed point). The symbol is now
$g(\mu)=p_k(\mu)$, and we are free to \emph{choose the polynomial}. The natural
demand is the minimax one: among all degree-$k$ polynomials with $p_k(0)=1$, make
the maximum of $\abs{p_k(\mu)}$ over the target window $[\lambda_{\min},
\lambda_{\max}]$ as small as possible.

\subsection{Why Chebyshev}

The minimax problem has a classical closed-form answer, and it is the reason these
polynomials are named after Chebyshev. The Chebyshev polynomials $T_k$ oscillate
between $-1$ and $+1$ on $[-1,1]$ and then climb away steeply outside that interval.
Affinely mapping the target window $[\lambda_{\min},\lambda_{\max}]$ onto $[-1,1]$
and normalizing so that $p_k(0)=1$ produces the unique polynomial that
\emph{equioscillates}---attains its small maximum amplitude many times---across the
window. No other degree-$k$ polynomial with a unit constant term is smaller across
the band. \Cref{fig:smoo-chebpoly} shows a degree-$4$ example: it is pinned small
over the target interval and rises sharply only at $\mu=0$, the smooth end we are
deliberately leaving to the coarse grid.

\begin{figure}[t]
\centering
\begin{tikzpicture}
\begin{axis}[
  width=12.6cm, height=6.0cm,
  xlabel={eigenvalue $\lambda$},
  ylabel={polynomial symbol $|p_k(\lambda)|$},
  xmin=0, xmax=1.05, ymin=0, ymax=1.05,
  xtick={0,0.2,0.4,0.6,0.8,1.0},
  extra x ticks={0.2,1.0}, extra x tick labels={$\lambda_{\min}$,$\lambda_{\max}$},
  extra x tick style={grid=major, grid style={simBlue!40, densely dotted},
    tick label style={text=simBlue, font=\scriptsize}},
  ytick={0,0.25,0.5,0.75,1},
  legend style={font=\scriptsize, at={(0.98,0.96)}, anchor=north east, draw=simGray!50},
  axis lines=left, tick style={simGray},
  samples=300, domain=0:1.05,
]
  % shaded target window [0.2, 1.0]
  \addplot[draw=none, fill=simBlue!10, forget plot] coordinates {(0.2,0)(1.0,0)(1.0,1.05)(0.2,1.05)} \closedcycle;
  \node[font=\scriptsize, simBlue, anchor=south] at (axis cs:0.6,0.86) {target window};
  % degree-1 (damped Jacobi-like affine) for contrast
  \addplot[simGray, thick, densely dashed] {abs(1 - 1.667*x)};
  % degree-4 chebyshev symbol: |T4((2 lam - lmin - lmax)/(lmax-lmin)) / T4(arg at 0)|
  % lmin=0.2 lmax=1.0 -> map m(lam) = (2 lam - 1.2)/0.8 ; T4(t)=8t^4-8t^2+1
  % normalize by T4 at lam=0: m(0) = -1.5, T4(-1.5)=8*5.0625-8*2.25+1=40.5-18+1=23.5
  \addplot[simBlue, very thick, samples=400] {abs((8*((2*x-1.2)/0.8)^4 - 8*((2*x-1.2)/0.8)^2 + 1)/23.5)};
  \legend{degree 1 (affine), degree 4 (Chebyshev)}
\end{axis}
\end{tikzpicture}
\caption[A Chebyshev polynomial shaped to the spectral window]{The degree-$4$
Chebyshev symbol $\abs{p_4(\lambda)}$ (blue) equioscillates across the target window
$[\lambda_{\min},\lambda_{\max}]$, holding the high-frequency error reduction below a
small uniform bound the whole way. A single affine step (gray dashed, the
damped-Jacobi symbol of \cref{sec:smoo-jacobi}) cannot stay small across the band.
Both rise toward $1$ only at $\lambda=0$, the smooth end left to the coarse grid.
Higher degree buys a smaller uniform bound across the window---at the cost of more
operator applies per sweep.}
\label{fig:smoo-chebpoly}
\end{figure}

The Chebyshev smoother has two further virtues that make it the multigrid smoother
of choice in Palace. It is \emph{matrix-free}: it applies $p_k(\mat{D}^{-1}\mat{A})$
by repeated operator applies, never forming the polynomial as a matrix. And, like
Jacobi, it is \emph{inner-product-free}: its coefficients are fixed closed forms of
the step index and the spectral bounds $[\lambda_{\min},\lambda_{\max}]$, computed
in advance without ever inspecting the iterate. No global reduction appears anywhere
in the sweep---a decisive advantage on a parallel or GPU machine, where inner
products force a synchronization that polynomial smoothing simply avoids. The only
global information it needs is an \emph{estimate} of $\lambda_{\max}$, computed once
at setup by a handful of power-iteration steps.

\begin{physicsbox}
The contrast with the conjugate-gradient method of \cref{ch:precond} is sharp and
deliberate. CG \emph{adapts} to the spectrum: it forms inner products at every
iteration and discovers the optimal polynomial as it goes. Chebyshev \emph{commits}
to a fixed polynomial in advance, paying for that with a spectral estimate but
buying back the inner products entirely. As a standalone solver CG wins. As a
multigrid smoother---applied thousands of times, on a machine where every inner
product is a synchronization---committing in advance is exactly the right trade.
\end{physicsbox}

\subsection{The three-term recurrence}

The polynomial is never written down or expanded. It is applied by a short
recurrence that threads three vectors: the current iterate \lstinline[style=l4style]|y|,
the residual \lstinline[style=l4style]|r|, and a search direction
\lstinline[style=l4style]|d|. One Chebyshev sweep of degree
\lstinline[style=l4style]|order| is:

\begin{lstlisting}[style=l4style]
-- one Chebyshev sweep: apply p_order(Dinv A) to the residual, accumulate into y.
-- coefficients sigma_d, sigma_r are closed forms of (k, lambda_max); no reductions.
chebyshev_sweep :: ChebOp[N] -> Tensor[N] -> Tensor[N] -> Tensor[N]
chebyshev_sweep op x y =
  let r0 = x `sub` (op.A `apply` y)        -- initial residual  r = x - A y
      d0 = scale (alpha0 op) (op.dinv `hadamard` r0)   -- first direction
  in  go 1 (y `add` d0) r0 d0              -- accumulate, then recur
  where
    go k y r d
      | k > op.order = y
      | otherwise =
          let r' = r `sub` (op.A `apply` d)            -- r -= A d
              d' = scale (sigma_d op k) d
                     `add` scale (sigma_r op k) (op.dinv `hadamard` r')
          in  go (k+1) (y `add` d') r' d'              -- y += d
\end{lstlisting}

Read the loop body as the three updates the chapter brief names: accumulate the
direction into the answer ($y \mathrel{+}= d$); update the residual by the
direction's operator image ($r \mathrel{-}= \mat{A}d$); and form the next direction
as a blend of the old direction and the diagonally-scaled residual
($d = \sigma_d\,d + \sigma_r\,(\mat{D}^{-1}r)$). Each iteration costs exactly one
operator apply ($\mat{A}d$) and a few elementwise vector operations. The
coefficients $\sigma_d, \sigma_r$ are pure functions of the step index $k$ and
$\lambda_{\max}$---for the fourth-kind variant Palace uses by default,
$\sigma_d = \tfrac{2k-1}{2k+3}$ and $\sigma_r = \tfrac{8k+4}{(2k+3)\lambda_{\max}}$.
\Cref{fig:smoo-chebdataflow} traces the dataflow of one step and marks the absence
of any reduction node.

\begin{figure}[t]
\centering
\begin{tikzpicture}[node distance=11mm and 16mm]
  \node[data] (r) {residual $\vr$};
  \node[opbox, right=of r] (Ad) {$\mat{A}\vd$ \\ \scriptsize(one apply)};
  \node[opbox, right=of Ad] (rup) {$\vr \mathrel{-}= \mat{A}\vd$};
  \node[opbox, below=10mm of Ad] (dscale) {$\mat{D}^{-1}\vr$ \\ \scriptsize(elementwise)};
  \node[opbox, right=of dscale] (dup) {$\vd = \sigma_d\vd + \sigma_r(\mat{D}^{-1}\vr)$};
  \node[product, right=18mm of rup] (yup) {$\vy \mathrel{+}= \vd$};
  \node[data, below=10mm of dup] (next) {next step};

  \draw[flow] (r) -- (Ad);
  \draw[flow] (Ad) -- (rup);
  \draw[flow] (rup) |- (dscale);
  \draw[flow] (dscale) -- (dup);
  \draw[flow] (dup) -- (yup);
  \draw[flow] (dup) -- (next);
  \draw[backflow] (rup.south) -- ++(0,-0.5) -| (next.north);

  % annotation: no reduction
  \node[cornertag, anchor=west, text=simGreen] at ($(yup.south east)+(0.1,-0.4)$)
    {\scriptsize no inner product};
  \node[cornertag, anchor=west, text=simGreen] at ($(yup.south east)+(0.1,-0.75)$)
    {\scriptsize no global sum};
\end{tikzpicture}
\caption[The Chebyshev recurrence is inner-product-free]{One step of the Chebyshev
three-term recurrence. The only global operation is the single operator apply
$\mat{A}\vd$; everything else is elementwise. Crucially there is \emph{no} inner
product and \emph{no} global reduction anywhere in the loop---the coefficients
$\sigma_d,\sigma_r$ are precomputed from the step index and $\lambda_{\max}$. This is
what makes Chebyshev the communication-light smoother of choice inside a
distributed multigrid V-cycle.}
\label{fig:smoo-chebdataflow}
\end{figure}

\begin{pitfall}
The recurrence is not a convenience; it is \emph{load-bearing numerics}. One could,
in principle, expand the polynomial into its monomial form
$p_k(\mat{D}^{-1}\mat{A})\vr = \sum_j c_j (\mat{D}^{-1}\mat{A})^{j}\vr$ and sum the
powers directly. \emph{Do not.} For the polynomial degrees in use the monomial sum
is numerically unstable---the high powers $(\mat{D}^{-1}\mat{A})^j$ overwhelm
floating-point precision and the coefficients $c_j$ alternate in sign and grow
enormous, so catastrophic cancellation destroys the result. The three-term
recurrence (a Clenshaw/Horner-style evaluation) computes the \emph{same} polynomial
in exact arithmetic but is stable in floating point. The recurrence form and the
monomial form are mathematically equal and computationally not interchangeable; the
sequentiality of the recurrence is part of the algorithm~\citep[\S2]{phillipsfischer2022}.
\end{pitfall}

\subsection{First-kind and fourth-kind variants}

Two flavors of Chebyshev smoother appear in practice, differing only in which
spectral window the minimax polynomial targets and hence in their coefficient
formulas. The \emph{first-kind} smoother targets a two-sided window
$[\lambda_{\min},\lambda_{\max}]$ and threads one extra scalar through its
recurrence; it needs an estimate of \emph{both} ends of the band. The
\emph{fourth-kind} smoother (the modern default, the one in the pseudocode above)
targets the one-sided window $[0,\lambda_{\max}]$ and needs only the top of the
spectrum---a single estimate, no lower bound to guess. Both produce a degree-$k$
polynomial small across the high-frequency band; the fourth-kind's one-sided design
makes it more robust when the spectral estimate is rough, which is why Palace adopts
it~\citep{phillipsfischer2022}. At the level of the synthesis surface the two are a
single operator whose coefficient generator is swapped---the variant is absorbed into
the closure, exactly as the damping mode was absorbed into Jacobi's
\lstinline[style=l4style]|dinv|.

\begin{workedexample}{A two-step Chebyshev smoother}{smoo-cheby}
Take a tiny system with $\mat{D}^{-1}\mat{A}$ already diagonal in the eigenvector
basis, with scaled eigenvalues spread across $[\lambda_{\min},\lambda_{\max}]=[0.2,1.0]$.
Build the degree-$2$ minimax polynomial $p_2$ with $p_2(0)=1$ that is smallest on
$[0.2,1.0]$. Map the window to $[-1,1]$ by $t(\lambda)=(2\lambda-1.2)/0.8$ and use
$T_2(t)=2t^2-1$. The normalized symbol is
\[
  p_2(\lambda) = \frac{T_2\bigl(t(\lambda)\bigr)}{T_2\bigl(t(0)\bigr)},
  \qquad t(0) = \frac{-1.2}{0.8} = -1.5,
  \qquad T_2(-1.5) = 2(2.25)-1 = 3.5,
\]
so $p_2(\lambda) = \bigl(2\,t(\lambda)^2 - 1\bigr)/3.5$. Evaluate the damping at three
frequencies:
\begin{center}
\setlength{\tabcolsep}{10pt}\renewcommand{\arraystretch}{1.15}
\begin{tabular}{@{}llll@{}}
\toprule
mode & $\lambda$ & $t(\lambda)$ & $\abs{p_2(\lambda)}$ \\
\midrule
smooth (left of window)   & $0.05$ & $-1.375$ & $0.795$ \\
mid-band                  & $0.6$  & $0$      & $0.286$ \\
high (top of window)      & $1.0$  & $1.0$    & $0.286$ \\
\bottomrule
\end{tabular}
\end{center}
The polynomial reduces every mode \emph{inside} the band $[0.2,1.0]$ to at most
$\approx 0.29$ of its amplitude in a single two-step sweep---compare damped Jacobi
from \cref{wex:smoo-jacobi}, which needed three sweeps to reach $0.33$ at the top
mode alone. The smooth mode at $\lambda=0.05$, outside the targeted band, is left at
$0.795$, untouched by design. Two operator applies have damped the entire
high-frequency half of the spectrum below $0.29$, uniformly. That uniformity across
the band---not just at one mode---is what the Chebyshev construction buys over a
single Jacobi step.
\end{workedexample}

\section{Hiptmair relaxation: smoothing in \texorpdfstring{$\Hcurl$}{H(curl)}}
\label{sec:smoo-hiptmair}

The smoothers so far assume one thing the eigenvector picture made explicit: that
the operator has \emph{no} large null space at the smooth end, so that ``high
frequency'' and ``oscillatory'' coincide and a diagonal relaxation can reach the
modes a coarse grid misses. For the scalar problems of electrostatics this holds.
For the vector curl--curl operator of magnetostatics, eigenmode, and driven
analyses, it fails---and the failure is not subtle.

\subsection{The gradient null space a plain smoother cannot see}

Recall from \cref{ch:bc} and the de Rham complex of \cref{ch:derham} that the
curl--curl operator $\Curl(\nu\Curl\,\cdot)$, discretized on Nédélec edge elements,
has an \emph{enormous null space}: every gradient field $\Grad\psi$ has zero curl,
so $\mat{A}\,\mat{G}\psi = 0$ for the discrete gradient $\mat{G}$ that maps scalar
potentials (nodal $\Hone$) into edge fields ($\Hcurl$). This is not a handful of
modes; it is an entire subspace, the discrete gradient range, with dimension on the
order of the number of mesh vertices.

A diagonal smoother is helpless against error living in that null space. Run the
damping analysis again: error along a null-space mode $\vphi$ has $\mat{A}\vphi=0$,
so the residual $\vr=-\mat{A}\ve$ has \emph{no component} along it, and the Jacobi
update \cref{eq:smoo-jac-update} does nothing---the damping factor is $g=1$ exactly,
for the entire gradient subspace. These are not high-frequency modes the smoother
merely damps weakly; they are modes it cannot touch \emph{at all}, and there are
vast numbers of them. \Cref{fig:smoo-nullspace} pictures the situation: the smoother
operates happily in the part of the space the curl ``sees,'' and is blind to the
gradient subspace beside it.

\begin{figure}[t]
\centering
\begin{tikzpicture}[scale=1.0]
  % the full H(curl) space as a big rounded rectangle
  \draw[thick, simInk, rounded corners=4pt, fill=simBgGray] (-4.4,-1.5) rectangle (4.4,1.7);
  \node[font=\small\bfseries, simInk, anchor=north west] at (-4.3,1.6) {$\Hcurl$ (edge dofs)};
  % the gradient null space subspace
  \draw[thick, simGreen, rounded corners=4pt, fill=simBgGreen] (-4.1,-1.25) rectangle (-0.4,1.0);
  \node[font=\footnotesize, simGreen, align=center] at (-2.25,0.55)
    {gradient null space\\$\operatorname{range}\mat{G}=\{\Grad\psi\}$\\\scriptsize($\mat{A}\,\Grad\psi = 0$)};
  % the rest: curl sees it
  \node[font=\footnotesize, simBlue, align=center] at (2.3,0.3)
    {curl-active part\\\scriptsize(smoother damps these)};
  % a "smoother" reaching only the right part
  \node[opbox, anchor=south] at (2.3,-1.35) {diagonal smoother};
  \draw[flow, simBlue] (2.3,-1.0) -- (2.3,-0.1);
  % blocked arrow into the null space
  \draw[->, simRust, thick] (-2.25,-1.0) -- (-2.25,-0.2);
  \node[font=\scriptsize, simRust] at (-2.25,-1.35) {$g=1$: no effect};
  \node[draw, cross out, simRust, thick, minimum size=4mm] at (-2.25,-0.55) {};
\end{tikzpicture}
\caption[The gradient null space a plain smoother misses]{The edge-dof space
$\Hcurl$ splits into the curl-active part, where a diagonal smoother damps error
normally, and the \emph{gradient null space} $\operatorname{range}\mat{G}$, where
$\mat{A}\,\Grad\psi=0$. On the null space the residual vanishes, the smoother's
damping factor is exactly $1$, and error simply persists. For an edge-element
problem this blind spot is a whole subspace---the smoother on its own cannot
converge. The cure is to relax in that subspace separately.}
\label{fig:smoo-nullspace}
\end{figure}

\begin{pitfall}
This is the specific reason edge-element ($\Hcurl$) problems need more than a
diagonal or even a Chebyshev smoother. The polynomial smoother of
\cref{sec:smoo-cheby} sharpens damping in the curl-active part but is just as blind
to the gradient null space as Jacobi---$p_k(0)=1$, and the entire null space sits at
$\lambda=0$. A multigrid built on a plain smoother will stall, its error trapped in
the gradient modes the coarse grid cannot drain because the smoother never made
them smooth. Magnetostatics, eigenmode, and driven analyses all hit this wall.
\end{pitfall}

\subsection{Relax in two spaces: the auxiliary-space idea}

The fix is due to Hiptmair, and in the de Rham setting of \cref{ch:derham} it is
beautifully natural: if the smoother cannot reach error in the gradient subspace
$\operatorname{range}\mat{G}$, then \emph{relax in that subspace too}, using the
discrete gradient $\mat{G}$ as the bridge. The discrete gradient maps the scalar
potential space (nodal $\Hone$) into the edge space ($\Hcurl$); its transpose
$\mat{G}^{\mathsf T}$ maps back. The combination
$\mat{A}_G = \mat{G}^{\mathsf T}\mat{A}\,\mat{G}$ is the curl--curl operator
\emph{seen from the scalar potential space}---a small, well-behaved scalar operator
with no null-space pathology of its own, on which an ordinary point smoother $B_G$
works perfectly.

One sweep of \emph{Hiptmair distributive relaxation} does two relaxations in
sequence (\cref{fig:smoo-hiptmair}):
\begin{align}
  \vy &\;\leftarrow\; \vy + B\,(\vx - \mat{A}\vy)
      && \text{primary (edge) relaxation,} \label{eq:smoo-hipt-prim}\\
  \vy &\;\leftarrow\; \vy + \mat{G}\,B_G\,\mat{G}^{\mathsf T}(\vx - \mat{A}\vy)
      && \text{auxiliary (gradient) correction.} \label{eq:smoo-hipt-aux}
\end{align}
The primary step \cref{eq:smoo-hipt-prim} is an ordinary smoother (Jacobi or
Chebyshev) on the edge operator $\mat{A}$---it handles the curl-active error. The
auxiliary step \cref{eq:smoo-hipt-aux} forms the residual, \emph{restricts} it to
the scalar space with $\mat{G}^{\mathsf T}$, relaxes there with $B_G$ on the
well-conditioned $\mat{A}_G$, and \emph{prolongs} the correction back into the edge
space with $\mat{G}$. Because $\mat{G}$ maps onto exactly the subspace the primary
smoother was blind to, the auxiliary step reaches precisely the error the first step
could not. Together they smooth the whole space. This is the discrete realization of
the de Rham auxiliary-space idea (Hiptmair--Xu) underlying the AMS
preconditioner~\citep{hiptmairxu2007}.

\begin{figure}[t]
\centering
\begin{tikzpicture}[node distance=9mm and 20mm]
  % primary (edge) space row
  \node[space, fill=simBgBlue, draw=simBlue] (Hcurl) {$\Hcurl$ \\ \scriptsize edge dofs};
  \node[opbox, right=of Hcurl, draw=simBlue] (B) {primary relax \\ $B$ on $\mat{A}$};
  \node[product, right=of B] (yout) {relaxed $\vy$};
  % auxiliary (scalar) space row, below
  \node[space, below=16mm of Hcurl] (H1) {$\Hone$ \\ \scriptsize nodal dofs};
  \node[opbox, right=of H1, draw=simGreen] (BG) {auxiliary relax \\ $B_G$ on $\mat{A}_G$};

  % residual into auxiliary space
  \draw[flow] (Hcurl) -- node[above,font=\scriptsize]{$\vx-\mat{A}\vy$} (B);
  \draw[flow] (B) -- (yout);
  % restrict G^T down
  \draw[flow, simGreen] (B.south) -- ++(0,-0.4) -| node[pos=0.25,above,font=\scriptsize]{restrict $\mat{G}^{\mathsf T}$} (H1.north);
  \draw[flow, simGreen] (H1) -- (BG);
  % prolong G up
  \draw[flow, simGreen] (BG.north) -- ++(0,0.4) -| node[pos=0.25,below,font=\scriptsize]{prolong $\mat{G}$} (yout.south);

  % labels
  \node[cornertag, anchor=west] at ($(Hcurl.north west)+(0,0.5)$) {\scriptsize primary space};
  \node[cornertag, anchor=west] at ($(H1.south west)+(0,-0.5)$) {\scriptsize auxiliary (gradient) space};
\end{tikzpicture}
\caption[Hiptmair primary + auxiliary relaxation]{One Hiptmair distributive
relaxation sweep. The primary smoother $B$ relaxes the edge operator $\mat{A}$ in
$\Hcurl$, damping the curl-active error. The residual is then restricted by the
discrete gradient transpose $\mat{G}^{\mathsf T}$ into the scalar space $\Hone$,
relaxed there by $B_G$ on the well-conditioned $\mat{A}_G=\mat{G}^{\mathsf T}\mat{A}\mat{G}$,
and the correction prolonged back by $\mat{G}$. The auxiliary leg reaches exactly the
gradient null space the primary leg could not. The two legs run in sequence: the
auxiliary residual is formed \emph{after} the primary update.}
\label{fig:smoo-hiptmair}
\end{figure}

In the synthesis surface the whole sweep is one constructed operator. It captures
the two operators $\mat{A}$ and $\mat{A}_G$, the discrete gradient $\mat{G}$, two
point smoothers $B$ and $B_G$ (each a Chebyshev closure), and a sweep count; its
apply is the two legs above, repeated:

\begin{lstlisting}[style=l4style]
-- Hiptmair distributive relaxation: relax in the primary edge space, then
-- correct through the discrete gradient G in the auxiliary scalar space.
hiptmair_smoother :: HiptmairOp[N,M] -> Tensor[N] -> Tensor[N] -> Tensor[N]
hiptmair_smoother op x y = iterate op.pc_it sweep y
  where
    sweep y =
      let y1 = y `add` (op.B `apply2` (x, y))          -- primary leg (edge)
          r  = x `sub` (op.A `apply` y1)               -- residual after primary
          xG = pin_ess op.ess_G (op.G `applyT` r)      -- restrict G^T r, pin Dirichlet dofs
          yG = op.B_G `apply` xG                        -- relax in scalar space
      in  y1 `add` (op.G `apply` yG)                    -- prolong G yG, accumulate
\end{lstlisting}

Two implementation points the brief calls out. The auxiliary residual is formed
\emph{after} the primary leg updates \lstinline[style=l4style]|y|, so the two legs
do not commute---the sweep is a \emph{multiplicative} (Gauss-Seidel-between-spaces)
composition, not an additive one; conflating them changes the operator. And the
restricted auxiliary residual \lstinline[style=l4style]|xG| has its essential
(Dirichlet) degrees of freedom pinned to zero before the auxiliary relaxation, the
same boundary-condition surgery \cref{ch:bc} performed on the primary system.

\begin{notationbox}
Palace's distributive smoother chooses Chebyshev for \emph{both} point smoothers
$B$ and $B_G$, never a Gauss--Seidel sweep. That is deliberate. A true distributive
relaxation would call for a triangular Gauss--Seidel solve in each space, but
Gauss--Seidel is a sequential forward/back substitution that does not parallelize
and runs poorly on a GPU. Replacing it with two inner-product-free Chebyshev sweeps
keeps the entire smoother communication-light, at the cost of an inexact---but
perfectly adequate---auxiliary relaxation in place of an exact auxiliary solve.
\end{notationbox}

\section{Smoothers as preconditioners and as multigrid components}
\label{sec:smoo-ascomponents}

Every smoother in this chapter is, structurally, an apply-function of the kind
\cref{ch:precond} defined: it consumes a vector and returns a vector, it is built
once at setup and applied many times, and---for the zero-initial-guess, single-sweep
case---it is \emph{linear} in its input, the exact shape a Krylov method wants in its
preconditioner slot. Indeed Palace installs the Jacobi smoother directly as a Krylov
preconditioner where a cheap diagonal scaling is all the problem needs.

But used as a standalone preconditioner a smoother is weak, and the eigenvector
picture says exactly why: it leaves the smooth, low-frequency error nearly
untouched. The damping factor returns to $1$ as $\lambda\to0$ for Jacobi
(\cref{fig:smoo-jacsymbol}), for Chebyshev (\cref{fig:smoo-chebpoly}, since
$p_k(0)=1$), and the gradient null space defeats both. A smoother makes the error
\emph{smooth}; it does not make it \emph{small}.

That leftover smooth error is precisely what a coarse grid is good at. \emph{Multigrid}
pairs the smoother with a coarse-grid correction: smooth on the fine grid to kill the
oscillatory error, project the smooth remainder down to a coarse grid where it looks
oscillatory and is cheap to resolve, solve there, and prolong the correction back.
\Cref{fig:smoo-vcycle} previews the structure---a V-cycle with a smoother bolted into
every level, descending and ascending. The smoother is the component this chapter
built; the recursion that surrounds it is the subject of \cref{ch:multigrid}.

\begin{figure}[t]
\centering
\begin{tikzpicture}[scale=1.0, every node/.style={font=\footnotesize}]
  % levels as horizontal bands at descending y, V shape
  \def\lw{1.6}
  % fine -> coarse going down-right, then back up
  \coordinate (f0) at (0,0);
  \coordinate (f1) at (1.3,-1.0);
  \coordinate (f2) at (2.6,-2.0);
  \coordinate (cc) at (3.9,-2.8);
  \coordinate (b2) at (5.2,-2.0);
  \coordinate (b1) at (6.5,-1.0);
  \coordinate (b0) at (7.8,0);
  % connecting V path
  \draw[flow, simGray] (f0)--(f1)--(f2)--(cc)--(b2)--(b1)--(b0);
  % smoother nodes on the way down
  \node[opbox, fill=simBgBlue, draw=simBlue] at (f0) {smooth};
  \node[opbox, fill=simBgBlue, draw=simBlue] at (f1) {smooth};
  \node[opbox, fill=simBgBlue, draw=simBlue] at (f2) {smooth};
  % coarse solve at the bottom
  \node[product, fill=simBgRust, draw=simRust] at (cc) {coarse \\ solve};
  % smoother nodes on the way up
  \node[opbox, fill=simBgBlue, draw=simBlue] at (b2) {smooth};
  \node[opbox, fill=simBgBlue, draw=simBlue] at (b1) {smooth};
  \node[product, fill=simBgGreen, draw=simGreen] at (b0) {result};
  % side labels
  \node[anchor=east, simGray] at ($(f0)+(-0.6,0)$) {finest grid};
  \node[anchor=north, simGray] at ($(cc)+(0,-0.55)$) {coarsest grid};
  \node[anchor=west, simGray, align=left] at ($(b0)+(0.5,0)$) {smooth = this\\chapter's\\smoother};
\end{tikzpicture}
\caption[Smoothers inside a multigrid V-cycle (preview)]{The multigrid V-cycle of
\cref{ch:multigrid}, previewed. Descending from the finest grid, a smoother runs at
each level to remove the oscillatory error; the smooth remainder is carried to the
coarsest grid where it is solved cheaply; ascending, the correction is prolonged
back with a smoother at each level again. Every blue node is a smoother of the kind
this chapter built---Jacobi, Chebyshev, or Hiptmair. The smoother is a poor solver
in isolation and the indispensable component here.}
\label{fig:smoo-vcycle}
\end{figure}

\begin{keyidea}
A smoother is built to be a multigrid \emph{component}, not a solver. Standalone it
stalls on smooth error; inside a V-cycle, paired with coarse-grid correction, that
weakness is exactly covered---the smoother handles the high frequencies the coarse
grid cannot resolve, the coarse grid handles the low frequencies the smoother cannot
damp. The division of labor is the design. This is why we measured smoothers by
their per-frequency damping factor rather than by how fast they solve: as a
component, damping the high band cheaply is the whole job.
\end{keyidea}

\summarysection
Writing the error in the eigenvector basis of $\mat{A}$ reduces a smoother to a
single curve: its per-frequency damping factor $g(\lambda)$. A \emph{smoother} is an
iteration whose $g$ is small at high frequency and near $1$ at low frequency, so it
cheaply removes oscillatory error while leaving smooth error---a deliberately
lopsided ``high-pass error filter'' that is a bad standalone solver but the heart of
multigrid. \emph{Damped Jacobi} scales the residual by $\omega\mat{D}^{-1}$; its
symbol is the affine $1-\omega\mu$, optimally damped at
$\omega^\star=2/(\mu_{\min}+\mu_{\max})$, and undamped Jacobi can fail to smooth at
all. \emph{Chebyshev} applies a fixed-degree minimax polynomial $p_k(\mat{D}^{-1}\mat{A})$
shaped to be uniformly small across a target spectral window; it is matrix-free,
inner-product-free, needs only a $\lambda_{\max}$ estimate, and must be evaluated by
its stable three-term recurrence rather than its unstable monomial expansion.
\emph{Hiptmair distributive relaxation} cures the gradient null space that defeats a
plain smoother on edge-element curl--curl problems, by relaxing in the primary edge
space and in the auxiliary scalar space bridged by the discrete gradient $\mat{G}$.
All three are apply-functions; standalone they are weak, but as the smoothing step of
the V-cycle in \cref{ch:multigrid} they are essential.

\furtherreading
The optimal damping analysis and the eigenvector view of stationary iterations are
standard; Saad~\citep[Ch.~4]{saad2003} is the canonical reference. The fourth-kind
Chebyshev smoother, its coefficient formulas, and the load-bearing stability of the
recurrence over the monomial form are developed by Phillips and
Fischer~\citep{phillipsfischer2022}, whose analysis Palace follows directly. The
auxiliary-space (AMS) framework for $\Hcurl$ and $\Hdiv$ problems, of which Hiptmair
distributive relaxation is the smoother realization, is laid out by Hiptmair and
Xu~\citep{hiptmairxu2007}; the de Rham structure it exploits is \cref{ch:derham},
and the V-cycle it slots into is \cref{ch:multigrid}.

\exercisessection
\begin{enumerate}[nosep]
  \item \textbf{Damping symbol.} For damped Jacobi on a model problem with scaled
  spectrum $\mu\in[0,2]$, write the symbol $g(\mu)=1-\omega\mu$. (a) For which
  $\omega$ does $\abs{g(\mu_{\max})}=\abs{g(2)}$ first exceed $1$? (b) Show that any
  $\omega>1$ leaves at least one mode amplified. (c) Conclude that
  $\omega\le1$ is necessary for convergence on this spectrum.

  \item \textbf{Optimal $\omega$ for the half-band.} Repeat the balance of
  \cref{eq:smoo-jac-optomega} for the target band $[\mu_{\mathrm{lo}},\mu_{\max}]
  =[1,2]$ and verify $\omega^\star=2/3$. Then compute the worst-case damping
  $\max_{\mu\in[1,2]}\abs{g(\mu)}$ at this $\omega^\star$. How many sweeps to drive
  the worst high-frequency mode below $10^{-3}$?

  \item \textbf{Why undamped fails.} Using the 1-D Laplacian eigenvalues
  $\lambda_k=4\sin^2(k\pi/2(n+1))$ from \cref{wex:smoo-jacobi}, compute the
  undamped ($\omega=1$) damping factor at the top mode $k=n$ as $n\to\infty$.
  What does it approach, and why does that make undamped Jacobi useless as a
  smoother even when it does not diverge?

  \item \textbf{Chebyshev versus Jacobi cost.} A degree-$k$ Chebyshev sweep costs
  $k$ operator applies; a damped-Jacobi sweep costs $1$. \Cref{wex:smoo-cheby}
  reached worst-case band damping $0.29$ in one degree-$2$ sweep ($2$ applies). How
  many damped-Jacobi sweeps (\cref{wex:smoo-jacobi}, worst-case $1/3$ per sweep)
  match that damping, and how many applies is that? Which wins on this band?

  \item \textbf{The monomial trap.} Explain in one paragraph why expanding
  $p_k(\mat{D}^{-1}\mat{A})$ into $\sum_j c_j(\mat{D}^{-1}\mat{A})^j$ and summing the
  powers is numerically unstable for moderate $k$, while the three-term recurrence of
  \cref{fig:smoo-chebdataflow} computing the same polynomial is stable. (Hint:
  consider the size of $\norm{(\mat{D}^{-1}\mat{A})^j}$ and the signs of the $c_j$.)

  \item \textbf{The gradient blind spot.} Let $\vphi=\mat{G}\psi$ be a gradient
  field, so $\mat{A}\vphi=0$. Show directly that a damped-Jacobi sweep
  \cref{eq:smoo-jac-update} starting from error $\ve=\vphi$ leaves $\ve$ unchanged.
  Then show that the auxiliary leg \cref{eq:smoo-hipt-aux} of the Hiptmair sweep
  \emph{does} reduce it, provided $B_G$ is a contraction on $\mat{A}_G$. (You may
  assume $\mat{G}^{\mathsf T}\mat{G}$ is invertible on the relevant subspace.)

  \item \textbf{Multiplicative, not additive.} The two Hiptmair legs
  \cref{eq:smoo-hipt-prim}--\cref{eq:smoo-hipt-aux} form the auxiliary residual from
  the \emph{updated} $\vy$. Write the single-sweep error-propagation operator for
  the multiplicative composition, and compare it to the (incorrect) additive version
  that uses the \emph{same} residual for both legs. On a mode in the gradient null
  space, which one converges and which one stalls? Why does the order matter?

  \item \textbf{Smoother as preconditioner.} Argue that for zero initial guess and a
  single sweep, the Chebyshev smoother $\vx\mapsto p_k(\mat{D}^{-1}\mat{A})\vx$ is a
  \emph{linear} operator (hence a valid Krylov preconditioner), while a Hiptmair
  sweep with nonzero initial guess is \emph{affine}, not linear. Which property does
  a multigrid V-cycle's smoothing step rely on, and which does an outer Krylov
  preconditioner slot require?
\end{enumerate}
